\documentclass[a4paper,12pt,twoside]{article}
\usepackage[left=1in,right=1in,top=1in,bottom=1in]{geometry}

\usepackage{xeCJK}
\setCJKmainfont{Songti TC}
\setCJKfamilyfont{kai}{Songti TC}
\newcommand{\kai}{\CJKfamily{kai}} % 用 \kai{使用楷書}

\usepackage{lmodern}
\usepackage{titling}
\usepackage{titlesec}
\usepackage{graphicx}
\usepackage{float}
\usepackage{amsmath, amsthm, amssymb, amsfonts, mathrsfs}
\usepackage{upgreek} % 直立的希臘字母

\usepackage[shortlabels]{enumitem}
\usepackage{setspace}
\renewcommand{\baselinestretch}{1.5}
\usepackage[colorlinks]{hyperref}
\usepackage{fancyvrb}
\usepackage{fancyhdr}
\usepackage{multicol}
\usepackage{caption}
\usepackage{subcaption}
\usepackage{listings}
\usepackage{bm}
\usepackage{bbm} % for \mathbbm{1} (indicator function)
\usepackage{centernot}
\usepackage{mdframed}
\usepackage{xltabular}
\usepackage{tcolorbox}
\tcbuselibrary{skins, breakable}
\usepackage{tikz}
\usetikzlibrary{patterns, arrows.meta, positioning, calc, decorations.pathmorphing, decorations.pathreplacing, decorations.markings, decorations.shapes, angles, quotes, 3d, matrix}
\usepackage[nameinlink, capitalise]{cleveref}
\usepackage{diagbox}
\usepackage{booktabs}
\usepackage{multirow,array}
\usepackage{tabularx}
\usepackage{mathtools}
\usepackage{indentfirst} % 首段 indent
\setlength{\parindent}{2em}
\usepackage{cancel} % 數學式子畫線槓掉
\newcommand\mcancel[2][red]{\setbox0=\hbox{$#2$}%
\rlap{\raisebox{.45\ht0}{\textcolor{#1}{\rule{\wd0}{1pt}}}}#2}
\setcounter{MaxMatrixCols}{20}
\usepackage{datetime2}
\usepackage{paracol}

\makeatletter
\def\tiangan#1{\expandafter\@tiangan\csname c@#1\endcsname}
\def\@tiangan#1{%
  \ifcase#1\or (甲)\or (乙)\or (丙)\or (丁)\or (戊)\or (己)\or (庚)\or (辛)\or (壬)\or (癸)\else(?)\fi
}
\AddEnumerateCounter{\tiangan}{\@tiangan}{(甲)}
\makeatother

% 支援 \wideparen 圓弧記號（若未載入 yhmath，可用此巨集穩定渲染）
\providecommand{\wideparen}[1]{\overset{\frown}{#1}}
\providecommand{\redtext}[1]{{\color{red} #1}}
\providecommand{\bluetext}[1]{{\color{blue} #1}}

\AtBeginDocument{%
 \setlength{\abovedisplayskip}{6pt}
 \setlength{\belowdisplayskip}{6pt}
 \setlength{\abovedisplayshortskip}{4pt}
 \setlength{\belowdisplayshortskip}{4pt}
}

\definecolor{codegreen}{rgb}{0,0.6,0}
\definecolor{codegray}{rgb}{0.5,0.5,0.5}
\definecolor{codepurple}{rgb}{0.58,0,0.82}
\definecolor{backcolour}{rgb}{0.95,0.95,0.92}

\lstdefinestyle{mystyle}{
 backgroundcolor=\color{backcolour},
 commentstyle=\color{codegreen},
 keywordstyle=\color{magenta},
 numberstyle=\tiny\color{codegray},
 stringstyle=\color{codepurple},
 basicstyle=\ttfamily\footnotesize,
 breakatwhitespace=false,
 breaklines=true,
 captionpos=b,
 keepspaces=true,
 numbers=left,
 numbersep=5pt,
 showspaces=false,
 showstringspaces=false,
 showtabs=false,
 tabsize=2
}
\lstset{style=mystyle}

\hypersetup{citecolor=blue, linkcolor=blue, urlcolor=blue}

\DeclareMathOperator{\supp}{supp}
\DeclareMathOperator{\skewness}{Skewness}
\DeclareMathOperator{\DU}{D-U}
\DeclareMathOperator{\CU}{C-U}
\DeclareMathOperator{\Geo}{Geo}
\DeclareMathOperator{\Pois}{Pois}
\DeclareMathOperator{\Gam}{Gam}
\DeclareMathOperator{\LN}{LogN}
\DeclareMathOperator{\cov}{Cov}
\DeclareMathOperator*{\argmax}{arg\,max}
\DeclareMathOperator*{\argmin}{arg\,min}
\DeclareMathOperator{\re}{Re}
\DeclareMathOperator{\im}{Im}

\newcommand{\N}{\mathbb{N}}
\newcommand{\Z}{\mathbb{Z}}
\newcommand{\Q}{\mathbb{Q}}
\newcommand{\R}{\mathbb{R}}
\newcommand{\C}{\mathbb{C}}
\newcommand{\G}{\mathcal{G}}
\newcommand{\mP}{\mathbb{P}}
\newcommand{\F}{\mathcal{F}}
\newcommand{\B}{\mathcal{B}}
\newcommand{\1}{\mathbbm{1}}
\newcommand{\Poi}{\text{Poi}}
\newcommand{\Var}{\text{Var}}
\newcommand{\as}{\text{a.s.}}
\newcommand{\vect}[1]{\boldsymbol{\mathbf{#1}}}
\newcommand{\paren}[1]{\left({#1}\right)}
\newcommand{\brk}[1]{\left[{#1}\right]}
\newcommand{\curl}[1]{\left\{{#1}\right\}}
\providecommand{\lambdaup}{\uplambda}
\renewcommand{\lambda}{\lambdaup}
\providecommand{\wideparen}[1]{\overset{\frown}{#1}}

\theoremstyle{definition} % 定理環境 style 改用正體
\newtheorem{theorem}{定理}
\newtheorem*{theorem*}{定理}
\newtheorem{definition}{定義}
\newtheorem*{definition*}{定義}
\newtheorem{example}{例子}
\newtheorem*{example*}{例子}
\newtheorem{exercise}{練習}
\newtheorem*{exercise*}{練習}
\newtheorem{corollary}{推論}
\newtheorem*{corollary*}{推論}
\newtheorem{property}{性質}
\newtheorem*{property*}{性質}
\newtheorem{prop}{命題}
\newtheorem*{prop*}{命題}
\newtheorem{proposition}[prop]{命題}
\newtheorem*{proposition*}{命題}
\newtheorem{lemma}{引理}
\newtheorem*{lemma*}{引理}
\newtheorem{problem}{問題}
\newtheorem*{problem*}{問題}
\newtheorem*{remark}{註}
\newtheorem{discussion}{論例}
\newtheorem*{discussion*}{論例}
\newtheorem{homework}{作業}
\newtheorem*{homework*}{作業}
\newtheorem{example-non}{例}

% 讓 \begin{solution} 與 \begin{solution*} 皆可用，並以 # 或 \square 結尾
\newtheorem*{solution}{【解】}
\AtEndEnvironment{solution}{\hfill$\#$}

% 每章（\section）自動重置定理、例題、論例、命題、公式等計數器，確保全冊與單章編譯編號完全一致
\makeatletter
\@addtoreset{theorem}{section}
\@addtoreset{definition}{section}
\@addtoreset{example}{section}
\@addtoreset{exercise}{section}
\@addtoreset{corollary}{section}
\@addtoreset{property}{section}
\@addtoreset{prop}{section}
\@addtoreset{lemma}{section}
\@addtoreset{problem}{section}
\@addtoreset{discussion}{section}
\@addtoreset{homework}{section}
\@addtoreset{example-non}{section}
\@addtoreset{equation}{section}
\makeatother

% 當原書該章各節（§ X-Y）獨立重新起算定理、例題、論例、公式編號時，於 \subsection 後呼叫 \resetsectcounters
\newcommand{\resetsectcounters}{%
  \setcounter{theorem}{0}%
  \setcounter{definition}{0}%
  \setcounter{example}{0}%
  \setcounter{exercise}{0}%
  \setcounter{corollary}{0}%
  \setcounter{property}{0}%
  \setcounter{prop}{0}%
  \setcounter{lemma}{0}%
  \setcounter{problem}{0}%
  \setcounter{discussion}{0}%
  \setcounter{homework}{0}%
  \setcounter{example-non}{0}%
  \setcounter{equation}{0}%
}

\renewcommand{\figurename}{圖}
\renewcommand{\tablename}{表}

% 統一勘誤表返回連結巨集：各單章主檔與全冊主檔皆於勘誤表標題處定義 \label{sec:current_errata}
\providecommand{\errataref}{\hyperref[sec:current_errata]{目錄前勘誤表}}

\DeclareMathOperator{\var}{Var}
\DeclareMathOperator{\corr}{corr}
\DeclareMathOperator{\rank}{rank}
\DeclareMathOperator{\MSE}{MSE}
\DeclareMathOperator{\SSE}{SSE}
\DeclareMathOperator{\SST}{SST}
\DeclareMathOperator{\SSR}{SSR}
\DeclareMathOperator{\sgn}{sign}
\DeclareMathOperator{\alt}{Alt}
\DeclareMathOperator{\adj}{adj}
\DeclareMathOperator{\curlOp}{curl}
\DeclareMathOperator{\spn}{span}
\DeclareMathOperator{\sse}{SSE}
\DeclareMathOperator{\tr}{tr}

\newcommand{\ep}{\varepsilon}
\newcommand{\p}{\partial}
\newcommand{\abs}[1]{\left|#1\right|}
\newcommand\norm[1]{\left\lVert#1\right\rVert}
\newcommand{\innerprod}[1]{\langle#1\rangle}
\newcommand{\ds}{\displaystyle}
\newcommand{\blank}{\rule[-1mm]{1.5cm}{0.1mm}}
\newcommand{\dint}{\displaystyle\int}
\renewcommand{\bar}{\overline}

\newtheoremstyle{remark}{}{}{}{}{\bfseries\sffamily}{.}{.5em}{#1}
\theoremstyle{remark}

\newtheorem*{remark*}{\textit{Remark}}
\newtheorem*{solution*}{解}

\AtEndEnvironment{solution*}{\hfill$\square$}
\newtheorem*{hint*}{Hint}


\fancypagestyle{firstpage}{
  \fancyhf{}
  \fancyhead[R]{\small《高中數學實驗教材》第 6 冊（社會組）}
  \fancyhead[L]{\small 第六冊\quad 索引}
  \fancyfoot[C]{\thepage}
}

\pagestyle{fancy}
\setlength{\headheight}{14pt}
\renewcommand{\headrulewidth}{0.5pt}
\renewcommand{\footrulewidth}{0pt}
\fancyhead{}
\fancyhead[LE]{\small 《高中數學實驗教材》第六冊（社會組）}
\fancyhead[RO]{\small 第六冊\quad 索引}
\fancyfoot[C]{\thepage}

\title{\LARGE\textbf{《高中數學實驗教材》第六冊（社會組）}\\[0.5em] \Large\textbf{第六冊\quad 索引}}
\author{\textbf{高中數學實驗教材編輯小組}\quad 編著}
\date{民國六十四年（1975）初版}

\begin{document}

\maketitle
\thispagestyle{firstpage}

% !TeX root = Appendix1_main.tex
% =========================================================================
% 《高中數學實驗教材》第六冊（社會組）
% 第六冊 索引
% =========================================================================
\phantomsection
\section*{第六冊\quad 索\qquad 引}\label{sec:b6s-app1}
\addcontentsline{toc}{section}{\texorpdfstring{第六冊\quad 索引}{第六冊 索引}}

\begin{multicols}{2}

\subsection*{一\quad 劃}
\noindent
$n$ 元組 \dotfill 1-1 乙\\
$n$ 維空間 $\mathbb{R}^n$ \dotfill 1-1 乙\\
$\vect{a}$ 的方向 \dotfill 1-1 乙\\
$n$ 維平行體 \dotfill 1-2 乙\\
（$n$ 階行列式） \dotfill 1-2 乙\\
一元 $n$ 次方程式 \dotfill 2-1\\
$n$ 元基本對稱多項式 \dotfill 2-2\\
Vandermonde 行列式 \dotfill 2-2\\
3 次方程式，4 次方程式的解法 \dotfill 2-4

\subsection*{二\quad 劃}
\noindent
$x$ 截距 \dotfill 2-2\\
$y$ 截距 \dotfill 2-2

\subsection*{三\quad 劃}
\noindent
子行列式 \dotfill 1-6\\
子式 \dotfill 1-6\\
下降 \dotfill 附錄-IV-4\\
上昇 \dotfill 附錄-IV-4\\
子空間 \dotfill 1-1 丙\\
$K$-重根 \dotfill 2-4

\subsection*{四\quad 劃}
\noindent
分量 \dotfill 1-1 乙\\
方向 \dotfill 1-1 乙\\
方陣（$n$ 階方陣） \dotfill 1-1 丁\\
方陣的乘法（線性映射的合成） \dotfill 1-7\\
方陣的加法（線性映射的加法） \dotfill 1-7\\
內積 \dotfill 1-8\\
方程式 \dotfill 2-1\\
方程式之一根 \dotfill 2-1\\
尤拉數 $e$ (Euler number) \dotfill 附錄-IV-3\\
不可微（分） \dotfill 附錄-IV-2\\
切線 \dotfill 附錄-IV-4

\subsection*{五\quad 劃}
\noindent
矛盾方程組 \dotfill 1-8\\
外積 \dotfill 附錄 I\\
代數的基本定理 \dotfill 2-2\\
平行變換 \dotfill 2-4\\
卡丹解法 \dotfill 2-5\\
可微（分） \dotfill 附錄-IV-2\\
史篤姆 (Sturm's) 定理 \dotfill 附錄 III

\subsection*{六\quad 劃}
\noindent
向量 \dotfill 1-1 乙\\
向量的加法 \dotfill 1-1 乙\\
列 \dotfill 1-1 丁\\
行 \dotfill 1-1 丁\\
列指標 \dotfill 1-1 丁\\
行指標 \dotfill 1-1 丁\\
行列式 \dotfill 1-2 乙\\
行列式的加法定理 \dotfill 1-2 乙\\
行列式的完全展開式 \dotfill 1-5\\
行列式的乘法定理 \dotfill 1-6\\
行列式的遞迴公式 \dotfill 1-6\\
多項方程式的求解 \dotfill 2-1\\
有向線段 \dotfill 1-1 乙\\
有向體積 \dotfill 1-2 乙\\
有向體積 \dotfill 1-7\\
有理係數方程式 \dotfill 2-1\\
共軛數 \dotfill 2-2\\
向上彎曲 \dotfill 附錄-IV-4\\
向下彎曲 \dotfill 附錄-IV-4\\
均值定理 \dotfill 附錄-IV-5

\subsection*{七\quad 劃}
\noindent
克萊瑪法則 \dotfill 1-8

\subsection*{八\quad 劃}
\noindent
直線 \dotfill 1-1 乙\\
始點 \dotfill 1-1 乙\\
法線 \dotfill 1-8\\
拉格蘭日 (Lagrange) 插入法 \dotfill 2-1\\
函數之圖示 \dotfill 2-2

\subsection*{九\quad 劃}
\noindent
重排 \dotfill 1-3\\
相依方程組 \dotfill 1-8\\
降冪方程式 \dotfill 2-2\\
軌跡 \dotfill 2-2\\
重根定理 \dotfill 2-4

\subsection*{十\quad 劃}
\noindent
原點 \dotfill 1-1 乙\\
座標系 \dotfill 1-1 丙\\
原素 \dotfill 1-1 丁\\
倒置 \dotfill 1-3\\
倒置數 \dotfill 1-3\\
剛性運動 \dotfill 1-7\\
高斯 (Gauss) \dotfill 2-2\\
根與係數的關係 \dotfill 2-3\\
特徵方程式 \dotfill 2-3\\
泰勒展開式 (Taylor's expansion) \dotfill 2-4\\
泰勒公式 \dotfill 附錄-IV-5

\subsection*{十一\quad 劃}
\noindent
終點 \dotfill 1-1 乙\\
第 $i$ 個單位向量 $\vect{e}_i$ \dotfill 1-1 乙\\
基底 \dotfill 1-1 丙\\
基準點 \dotfill 1-1 丙\\
第 $i$ 個列向量 \dotfill 1-1 丁\\
第 $j$ 個行向量 \dotfill 1-1 丁\\
符號函數 \dotfill 1-3\\
插入法 \dotfill 2-1\\
第一階導函數 \dotfill 2-4

\subsection*{十二\quad 劃}
\noindent
超平面 \dotfill 1-8\\
極大極小 \dotfill 2-2\\
極小值 $f(x_0)$ \dotfill 2-4\\
極大值 $f(x_0)$ \dotfill 2-4\\
極限 \dotfill 附錄-IV-1\\
極值 \dotfill 附錄-IV-4\\
極大值 \dotfill 附錄-IV-4\\
極小值 \dotfill 附錄-IV-4\\
結式 \dotfill 附錄 III

\subsection*{十三\quad 劃}
\noindent
解方程式 \dotfill 2-1\\
微分 \dotfill 附錄-IV-2

\subsection*{十四\quad 劃}
\noindent
端點 \dotfill 1-1 乙\\
複係數方程式 \dotfill 2-1\\
對稱性 \dotfill 2-2\\
對稱式 \dotfill 2-2\\
對稱式基本定理 \dotfill 2-3\\
對稱多項式 \dotfill 2-3\\
對稱式 \dotfill 2-3

\subsection*{十五\quad 劃}
\noindent
實係數方程式 \dotfill 2-1\\
餘因式 \dotfill 1-6

\subsection*{十六\quad 劃}
\noindent
線段 \dotfill 1-1 乙\\
線性組合 \dotfill 1-1 丙\\
線性相關 \dotfill 1-1 丙\\
線性無關 \dotfill 1-1 丙\\
線性映射 \dotfill 1-1 丁\\
線性空間 \dotfill 1-8\\
導函數 \dotfill 2-4\\
導數 \dotfill 附錄-IV-2

\subsection*{二十三\quad 劃}
\noindent
體積 \dotfill 1-7

\end{multicols}


\end{document}
